%% OVERFORE-Subdocumento6-Anexos.tex — generado desde OVERFORE-Subdocumento6-Anexos.docx
%% Formato: clase overfore, modo anexo. El formato lo pone la clase; aquí sólo va el contenido.
\documentclass[anexo]{../../comun/overfore}
\numsubdoc{6}
\informe{Informe Preparatorio N° 2}
\fecha{12 de octubre de 2026}
\externos{S6}
\begin{document}

\anexo{Cálculo de capacidad temporal y programación de sprints}
Este anexo detalla la carga de cada equipo de desarrollo y el cálculo de capacidad temporal con que se programan los sprints de cada etapa, que el cuerpo del Subdocumento 6 presenta en la sección 6.2.1. Las referencias a figuras y tablas con numeración del capítulo 6 remiten al cuerpo del Subdocumento 6.

\seccion{Carga de los equipos de desarrollo}
Con la organización por cadena de valor de la Figura \refext{S6}{fig:equipos-desarrollo-modulos-construyen} del Subdocumento 6, el equipo de cadena del mineral concentra 52 requerimientos funcionales de la etapa 1, los 51 de los tres módulos de la cadena más el de tableros de turno y CIO, y cuenta con la mayor dotación en esa etapa. El líder de integración responde además por el integrador de sistemas y por la coordinación técnica con los dueños de cada sistema existente. El paquete 1.4.3.7 suma 768 HH entre los meses 12 y 17 para corregir lo que cada equipo construyó en la etapa 1 durante su marcha blanca.

\seccion{Cálculo de capacidad temporal y programación de sprints}
Con sprints de dos semanas, nueve meses × 4,33 semanas por mes / 2 semanas por sprint equivalen a 19,49 sprints, y se programan 19 completos para la etapa 1, entre los meses 3 y 11. Para cinco meses, el mismo cálculo equivale a 10,83 sprints, y se programan 10 completos para la etapa 2, entre los meses 13 y 17. Las fracciones restantes absorben coordinación y cierres dentro del plazo contractual. Los primeros sprints de la etapa 2 trabajan sobre historias ya definidas mientras se cierra su línea base de alcance, que se aprueba en el mes 14.

\end{document}
